\scanpage{31}
\mspage{14}
On peut prendre aussi comme seuls gén. $\rho,\sigma_0,\tau_0$ (dont se déduisent $\sigma_1,\tau_1$, \inserted{$\sigma_\infty,\tau_\infty$} comme $\rho\sigma_0\rho^{-1},\rho\tau_0\rho^{-1},\rho^2\sigma_0\rho^{-2},\rho^2\tau_0\rho^{-2}$) et les relations
\[
\left\{\begin{aligned}
\unsure{\text{i)}|\text{a)}}\quad&\sigma_0^2=\tau_0^2=1\quad\tau_0\sigma_0=\sigma_0\tau_0\quad\text{(i.e. }(\overbrace{\sigma_0\tau_0}^{\tilde\sigma_0})^2=1\text{)}\\
\unsure{\text{ii)}|\text{b)}}\quad&\rho^3=1\\
&\sigma_0\underbrace{\rho\sigma_0^{-1}\rho^{-1}}_{\sigma_1}=\underbrace{\rho^2\tau_0\rho^{-2}}_{\tau_\infty}\underbrace{\rho\tau_0\rho^{-1}}_{\tau_1}\rho\\
&\text{i.e. }\sigma_0\rho\sigma_0^{-1}\rho^{-1}=\rho^{-1}\tau_0\rho\tau_0\\
&\text{ou enfin }\rho\sigma_0\rho^{-1}\sigma_0=\tau_0\rho^{-1}\tau_0\rho\quad\text{soit}\\
&\text{soit}\quad[\rho,\sigma_0][\rho^{-1},\tau_0]=1
\end{aligned}\right.
\]
\margin{\deleted{\unsure{i.e.}} ou encore $\deleted{\operatorname{int}(\rho)\sigma_0}\quad[\rho,\sigma_0]=[\tau_0,\rho^{-1}]$}

C'est donc un quotient du produit libre $(\mathbb Z/2\mathbb Z\times\mathbb Z/2\mathbb Z)*\mathbb Z/3\mathbb Z$, par le sous-groupe invariant décrit par la relation $[\rho,\sigma_0][\rho^{-1},\tau_0]=1$.

\uline{NB} On trouve une symétrie dans le groupe, qui échange $\sigma_0$ et $\tau_0$, et échange $\rho$ et $\rho^{-1}$. [donc qui échange aussi $\sigma_1$ et $\tau_\infty$, $\sigma_\infty$ et $\tau_1$] --- mais cette symétrie n'est pas compatible avec la structure d'extension, comme on voit de suite (en en passant au quotient $\unsure{\mathfrak S_3|\mathcal S_3}\times\mathbb Z/2$, elle devrait conserver le centre!).
\mspage{30}
